\section{Геометрические приложения интеграла Римана}
\subsection{Площади}
\begin{definition}
	\textit{Площадь} - мера Жордана в $\R^2$. Измеримые по Жордану множества называются \textit{квадрируемыми}.
\end{definition}
\begin{reminder}
	$L(P,f)$ - нижняя сумма Дарбу.
\end{reminder}
\begin{theorem}\label{25_th1}
	Если $f$ - неотрицательная, $f\in\mc R[a,b]$ функция, то криволинейная трапеция $G=\{(x,y):a\le x\le b,\ 0\le y\le f(x)\}$ квадрируема, причём:
	\[\mu_J(G)=\int_a^bf(x)dx\]
\end{theorem}
\begin{proof}
	Знаем из (\ref{7_th2}):
	\[fa\eps>0)(\ex P)\ U(P,f)-L(P,f)<\eps\]
	При этом:
	\[L(P,f)\le\int_a^bf(x)dx\le U(P,f),\ L(P,f)\le \mu_*^J(G)\le\mj(g)\le U(P,f)\]
	Тогда имеем:
	\[\mu_J(G)=\int_a^bf(x)dx\]
	Получили требуемое.
\end{proof}
\begin{definition}
	Площадь в полярных координатах.\\
	$G:=\{(\rho,\phi):\al\le\phi\le\beta,\ \rho\in[0,\rho(\phi)]\},\ \beta-\al\le2\pi$ при $x=\rho\cdot\cos\phi,\ y=\rho\cdot\sin\phi$
\end{definition}
\begin{exercise}
	Найдите площадь кругового сектора $\phi_i\le\phi\le\phi_{i+1},\ \rho\in[0,r]$ с помощью \nf{(\ref{25_th1})}.
\end{exercise}
\begin{theorem}
	Если $\rho$ - неотрицательная и непрерывная на $[\al,\beta],\ \beta-\al\le2\pi$ функция, то $G$, заданная с помощью $\rho(\phi)$ квадрируема, причём:
	\[\mu_J(G)=\frac12\int_\al^\beta\rho^2(\phi)d\phi\]
\end{theorem}
\begin{proof}
	Площадь кругового сектора $\phi_i\le\phi\le\phi_{i+1},\ \rho\in[0,r]$: \[\frac12r^2\tr\phi_i\text{, где }\tr\phi_i=\phi_{i+1}-\phi_i\]
	Возьмём $P:\al=\phi_0<\phi_1<\dots<\phi_n=\beta$. Рассмотрим $\ds L\l(P,\ \frac12\rho^2\r)$.\\
	На каждом рассматриваемом промежутке из $P$ точная нижняя грань будет соответствовать радиусу сектора круга с минимальным радиусом, который попадает в текущий промежуток, тогда:
	\[L\l(P,\ \frac12\rho^2\r)\le\mu^J_*(G)\le\mj(G)\le U\l(P,\ \frac12\rho^2\r)\]
	Аналогично с (\ref{25_th1}). Получили требуемое.
\end{proof}
\begin{definition}
	Пусть $\gamma$ - гладкая кривая в $\R^2$ или $\R^3$ с параметризацией $\v r(t),\ t\in[t_0,T]$. Для длины дуги кривой $\gamma$ справедлива формула:
	\[S'(t)=|\v r'(t)|\]
	Следовательно, \textit{длина кривой} $\gamma$ равна:
	\[\int_{t_0}^Ts'(t)dt=\int_{t_0}^T|\v r'(t)|dt\]
\end{definition}
\begin{note}
	Длина кривой определялась как точная верхняя грань вписанных в неё ломаных. Тогда, если рассуждать по аналогии, площадь поверхности - точная верхняя грань площадей всех вписанных многогранников, однако это неправильно, покажем это в следующем примере.
\end{note}
\begin{example}\nf{(Сапог Шварца)}\\
	Рассмотрим цилиндр с радиусом основания 1 и высотой 1, разделим высоту на $n$ равных частей, каждую из получившихся окружностей на $m$ равных частей. \\
	Как делить? Пусть первая окружность разделена на равные дуги, каждая углом $\frac{2\pi}m$. На второй окружности точки поставлены над серединами дуг первой окружности. Каждая следующая окружность под нечётным номером разделена также, как и первая, каждая следующая под четным номером - так же как и вторая.\\
	Теперь соединим эти точки у каждых двух соседних окружностей, получим треугольники. Найдём площадь одного такого.\\
	Пусть точки $A$ и $B$ - соседние на окружности $p$, $C$ лежит на окружности над $p$. \\
	Ищем $S_{\tr ABC}$.\\
	Пусть $O$ - центр $p$, $D$ - проекция $C$ на $p$, $H$ - пересечение $AB$ и $OD$. Тогда имеем:
	\[CD=\frac1n,\ CH=\sqrt{\frac1{n^2}+(1-\cos\pi m)^2},\ S_{\tr ABC}=\sin\frac\pi m\cdot\sqrt{\frac1{n^2}+(1-\cos\pi m)^2}\]
	Тогда площади поверхности:
	\[S_{\text{полн.}}=2mn\sin\frac\pi m\cdot\sqrt{\frac1{n^2}+(1-\cos\pi m)^2}\]
	Заметим:
	\[S_{\text{полн.}}\ra2\pi\cdot\sqrt{1+\frac{\pi^4p^2}4}\text{ при }n\ra\i \text{ и }\frac n{m^2}\ra p\]
	При выборе разных $p$ можно получить разные предельные значения площади, поэтому естественное определение площади поверхности не подходит.
\end{example}
\begin{anote}
	Треугольники из примера выше можно представить как украшение у барабана: \href{https://eliopt.ru/assets/images/Karlsbach/KB%202022/08870.jpg}{вот пример}.
\end{anote}
\begin{definition}
	\textit{Площадью поверхности тела вращения} $\mc D$ (из \ref{25_th3}) назовём предел площадей поверхностей фигур, полученных в результате вращения ломаных, вписанных в график функции $y=f(x)$, когда длина максимума проекции каждого звена на $O_x$ стремится к 0.
\end{definition}
\begin{theorem}
	Если $f$ - неотрицательная, непрерывно дифференцируемая на $[a,b]$ функция, то площади поверхности тела $\mc D$ \nf{(из \ref{25_th3})} равна: \[S=\ds2\pi\int_a^bf(x)\sqrt{1+(f'(x))^2}dx\]
\end{theorem}
\begin{proof}
	Площадь усечённого конуса, при центре большей окружности на $x_{i+1}$, а меньшей на $x_i$:
	\[\pi\Big(f(x_i)+f(x_{i+1})\Big)\sqrt{\Big(x_{i+1}-x_i\Big)^2+\Big(f(x_{i+1})-f(x_i)\Big)^2}\]
	Требуется доказать:
	\[\lim_{\tr(P)\ra0}\pi\sum_{i=1}^n\Big(f(x_i)+f(x_{i+1})\Big)\sqrt{\Big(x_{i+1}-x_i\Big)^2+\Big(f(x_{i+1})-f(x_i)\Big)^2}=S\]
\end{proof}
\subsection{Объёмы}
\begin{definition}
	\textit{Объём} - мера Жордана в $\R^3$. Измеримые по Жордану множества называются \textit{кубируемыми}.
\end{definition}
\begin{example}
	Цилиндрическое тело: $\mc D=D\times[a,b]$, где $D$ - квадрируемая фигура в $\R^2$.\\
	Докажем, что $\mc D$ кубируемо, причём:
	\[\mu_J(\mc D)=\mu_J(\mc D)\cdot(b-a)\]
	$D$ - квадрируемо, следовательно:
	\[(\fa\eps>0)(\ex\text{ элементарные }E_1,E_2:E_1\subset D\subset E_2)\ |E_2\sm E_1|<\frac\eps{b-a}\]
	Положим:
	\[\mc E_1:=E_1\times[a,b],\ \mc E_2:=E_2\times[a,b],\ \mc E_1\subset \mc D\subset \mc E_2\]
	\[|\mc E_2\sm\mc E_1|=|E_2\sm E_1|\cdot(b-a)<\eps\]
	\[|\mc E_1|=|E_1|\cdot(b-a)\le\mu_J(\mc D)\le|E_2|\cdot(b-a)=|\mc E_2|\]
	\[|E_1|\cdot(b-a)\le\mu_J(\mc D)\cdot(b-a)\le|E_2|\cdot(b-a)\]
	Получили требуемое.
\end{example}
\begin{theorem}\nf{(Объём тела вращения)}\label{25_th3}\\
	Если $f$ - неотрицательная и непрерывная на $[a,b]$ функция, то фигура, полученная в результате вращения криволинейной трапеции $\mc D=\{(x,y):a\le x\le b,\ 0\le y\le f(x)\}$ вокруг оси $O_x$ кубируема, причём:
	\[\mu_J(\mc D)=\pi\int_a^bf^2(x)dx\]
\end{theorem}
\begin{proof}
	Пусть $P:a=x_0<x_1<\dots<x_n=b$, рассмотрим $L(P,\ \pi f^2)$:
	\[L(P,\ \pi f^2)=\sum_{k=1}^n\pi m^2_k\tr x_k\le\mu_*^J(\mc D)\le\mj(\mc D)\le U(P,\ \pi f^2),\ m_k=\min_{x\in[x_{k-1}, x_k]}f(x)\]
\end{proof}